\documentclass[11pt,a4paper]{article}

\usepackage[a4paper,margin=14mm]{geometry}
\usepackage{fontspec}
\setmainfont{Arial}
\usepackage{amsmath,amssymb}
\usepackage{enumitem}
\usepackage{xcolor}
\usepackage{tikz}
\usepackage[most]{tcolorbox}
\usepackage{ragged2e}
\usepackage{fancyhdr}
\usepackage{eso-pic}
\usepackage{array}
\usepackage{tabularx}
\usepackage[hidelinks]{hyperref}

\definecolor{BrandPrimary}{HTML}{A02050}   % warmer, slightly lighter maroon — less heavy on the eye
\definecolor{BrandSecondary}{HTML}{4A7FD4} % softer muted blue — less corporate, more educational
\definecolor{BrandGold}{HTML}{C9A56A}
\definecolor{BrandSoft}{HTML}{F8F2E7}
\definecolor{Ink}{HTML}{2D3748}            % slightly warmer charcoal for body text
\definecolor{Muted}{HTML}{8794A3}          % lighter muted — gentler auxiliary text
\definecolor{Line}{HTML}{E2E8F0}           % warm-tinted light border

\pagestyle{fancy}
\fancyhf{}
\fancyfoot[L]{\footnotesize\color{Muted}\href{https://www.25maths.com}{www.25maths.com}}
\newcommand{\SetFooterLabel}[1]{%
  \fancyfoot[R]{\footnotesize\color{Muted}#1}%
}
\SetFooterLabel{Student Sheet 1/2}
\renewcommand{\headrulewidth}{0pt}

\newcommand{\WorksheetTitle}{Number Arithmetic: BIDMAS, HCF, LCM}
\newcommand{\WorksheetSubtitle}{Kahoot Extension Worksheet \;\textbar\; Topic ID: 01 \;\textbar\; Difficulty: Auto}
\newcommand{\BoardLabel}{CIE 0580 / Edexcel 4MA1}

\newcommand{\BrandFrame}{%
  \AddToShipoutPictureBG{%
    \begin{tikzpicture}[remember picture,overlay]
      \draw[BrandPrimary!70!white,line width=0.9pt,rounded corners=2.0mm] ([xshift=8mm,yshift=-8mm]current page.north west) rectangle ([xshift=-8mm,yshift=8mm]current page.south east);
    \end{tikzpicture}%
  }%
}

\newtcolorbox{topbandbox}{
  colback=BrandSoft,
  colframe=BrandGold,
  arc=2.2mm,
  boxrule=0.5pt,
  left=2.2mm,right=2.2mm,top=1.7mm,bottom=1.7mm
}

\newtcolorbox{focusbox}{
  colback=BrandSecondary!4!white, % very gentle blue tint — softer than pure blue!3
  colframe=BrandSecondary!55!white, % reduced border saturation for a calmer feel
  arc=2.0mm,   % rounder corners — friendlier visual
  boxrule=0.5pt,
  left=2.2mm,right=2.2mm,top=1.6mm,bottom=1.6mm
}

\newtcolorbox{notesbox}{
  colback=BrandSoft!50!white,  % slightly softer cream
  colframe=BrandGold!70!white, % lighter gold border — less heavy
  arc=2.0mm,   % match focusbox rounding
  boxrule=0.4pt,
  left=2.2mm,right=2.2mm,top=1.4mm,bottom=1.4mm
}

\newtcolorbox{fieldbox}[1]{
  colback=white,
  colframe=Line,
  arc=1.8mm,       % rounder to match overall softness
  boxrule=0.4pt,
  left=1.8mm,right=1.8mm,top=1.1mm,bottom=1.1mm,
  height=14.5mm,
  valign=center,
  halign=left,
  fonttitle=\footnotesize\bfseries,
  coltitle=Muted,
  title={#1}
}

\newcommand{\ansline}{\textcolor{Line!80!white}{\rule{\linewidth}{0.24pt}}} % thinner, slightly faded lines
\newcommand{\answerlines}[1]{%
  \vspace{1.4mm}
  \begingroup
  \setlength{\parskip}{2.4mm}  % more vertical space between lines — less cramped
  \setlength{\parindent}{0pt}
  \foreach \n in {1,...,#1}{\ansline\par}
  \endgroup
}

\newtcolorbox{qbox}[2][]{
  enhanced,
  breakable,
  width=\linewidth,
  colback=white,
  colframe=Line,
  arc=2.0mm,          % rounder corners — consistent with focus/notes boxes
  boxrule=0.4pt,      % slightly thinner border — lighter visual weight
  left=2.4mm,right=2.4mm,top=1.8mm,bottom=2.0mm,  % a touch more breathing room
  title={\textbf{#2}},
  coltitle=BrandPrimary,
  fonttitle=\bfseries\normalsize,
  before upper={\RaggedRight\setlength{\parindent}{0pt}\hyphenpenalty=10000\exhyphenpenalty=10000\emergencystretch=2em},
  #1
}

\begin{document}
\color{Ink}
\BrandFrame

\begin{topbandbox}
\begin{minipage}[t]{0.56\textwidth}
\textbf{\color{BrandPrimary}25Maths Worksheet}
\end{minipage}
\hfill
\begin{minipage}[t]{0.4\textwidth}
\raggedleft\footnotesize\color{Muted}
Board: \BoardLabel\\
Format: Printable A4 PDF
\end{minipage}
\end{topbandbox}

\vspace{2mm}
\begin{center}
{\LARGE\bfseries\color{BrandPrimary}\WorksheetTitle}\par
\vspace{0.6mm}
{\small\bfseries\color{Muted}\WorksheetSubtitle}
\end{center}

\vspace{3mm}
\noindent
\begin{minipage}[t]{0.32\textwidth}
\begin{fieldbox}{Student Name}\vspace*{2.8mm}\end{fieldbox}
\end{minipage}\hfill
\begin{minipage}[t]{0.32\textwidth}
\begin{fieldbox}{Class / Group}\vspace*{2.8mm}\end{fieldbox}
\end{minipage}\hfill
\begin{minipage}[t]{0.32\textwidth}
\begin{fieldbox}{Date}\vspace*{2.8mm}\end{fieldbox}
\end{minipage}

\vspace{2mm}
\begin{focusbox}
\textbf{Syllabus focus:} Apply order of operations, calculate HCF/LCM, and convert between ordinary form and standard form using exam-style wording.
\end{focusbox}

\vspace{1.5mm}
\begin{notesbox}
\textbf{Exam reminder:} Powers first, then multiplication/division, then addition/subtraction. In word problems, decide whether you need \textbf{HCF} (grouping) or \textbf{LCM} (cycle timing).
\end{notesbox}

\vspace{2mm}
\noindent{\large\bfseries Practice (10 Questions)}

\begin{qbox}{Q1 (BIDMAS)}
Evaluate ~$14-3\times2$~.
\answerlines{2}
\end{qbox}

\begin{qbox}{Q2 (Brackets)}
Evaluate ~$\left(\frac{18}{3}\right)(5-2)$~.
\answerlines{2}
\end{qbox}

\begin{qbox}{Q3 (HCF)}
Find the highest common factor of ~$24$~ and ~$36$~.
\answerlines{2}
\end{qbox}

\begin{qbox}{Q4 (LCM)}
Find the lowest common multiple of ~$8$~ and ~$12$~.
\answerlines{2}
\end{qbox}

\begin{qbox}{Q5 (Standard Form)}
Write ~$0.0032$~ in standard form.
\answerlines{2}
\end{qbox}

\begin{qbox}{Q6 (Powers of 10)}
Write ~$7.5\times10^2$~ as an ordinary number.
\answerlines{2}
\end{qbox}

\begin{qbox}{Q7 (Rounding)}
Round ~$12.467$~ to ~$2$~ decimal places.
\answerlines{2}
\end{qbox}

\begin{qbox}{Q8 (Integers)}
Evaluate ~$20-(-4)$~.
\answerlines{2}
\end{qbox}

\begin{qbox}{Q9 (Estimation)}
Estimate ~$\dfrac{49.8}{10.1}$~.
\answerlines{2}
\end{qbox}

\begin{qbox}{Q10 (LCM in Context)}
Two buses leave every ~$15$~ min and ~$20$~ min. They leave together at 09:00. When is the next same departure time?
\answerlines{3}
\end{qbox}

\vspace{1.5mm}
\noindent{\large\bfseries Quick Answer Grid}

\vspace{1.5mm}
\setlength{\arrayrulewidth}{0.35pt}
\setlength{\tabcolsep}{0pt}
\renewcommand{\arraystretch}{1.55}
\noindent
\begin{tabularx}{\linewidth}{|*{10}{>{\centering\arraybackslash}X|}}
\hline
Q1 & Q2 & Q3 & Q4 & Q5 & Q6 & Q7 & Q8 & Q9 & Q10 \\
\hline
\ \ & \ \ & \ \ & \ \ & \ \ & \ \ & \ \ & \ \ & \ \ & \ \ \\
\hline
\end{tabularx}

\end{document}
